% Source: https://artofproblemsolving.com/community/c3h1734624_2019_mathcounts_handbook_warmup_3_8_and_10
% Stable identity: aops-topic-1734624-q02
From each of the $17$ vertices, diagonals go to every vertex except itself and its two neighbours, that is to $14$ vertices. This gives $17\cdot14$ diagonal ends, and each diagonal has two ends, so the number of diagonals is
\[
\dfrac{17\cdot14}{2}=\boxed{119}.
\]
(The number $14$ is the number of diagonals from a single vertex, not the total; the figure shows these $14$ diagonals from one vertex $V$ of a $17$-gon.)

\begin{asy}
usepackage("amsmath");
size(5cm);
defaultpen(fontsize(10pt));
pair[] P;
for(int i=0;i<17;++i) P.push(dir(90+360*i/17));
draw(P[0]--P[1]--P[2]--P[3]--P[4]--P[5]--P[6]--P[7]--P[8]--P[9]--P[10]--P[11]--P[12]--P[13]--P[14]--P[15]--P[16]--cycle,linewidth(0.9));
for(int i=2;i<16;++i) draw(P[0]--P[i],blue+linewidth(0.5));
for(int i=0;i<17;++i) dot(P[i],linewidth(3));
label("$V$",P[0],N);
label("\mbox{}",(min(currentpicture,user=true)+max(currentpicture,user=true))/2,black); // keeps the text after the figure black when LaTeX typesets the labels inline
\end{asy}
